% ---- Analysis ---------------------------------------------------------------
% Write qualitative analysis after training. Placeholders for now.

\subsection{Qualitative Examples}

\todo{Insert a figure grid: 3--4 rows of (prompt, generated SVG, ground-truth icon)
showing representative successes across icon styles.
Row 1: simple monochrome outline icons (e.g., arrows, UI actions)
Row 2: multicolor icons or icons with badges/decorators
Row 3: abstract/metaphorical icons (e.g., "speed", "security")
Row 4: multi-element compositions}

\begin{figure}[h]
    \centering
    \todo{Insert figure: qualitative\_examples.pdf}
    \caption{Representative \method{} outputs alongside their prompts
    and the closest ground-truth icons from the test split.}
    \label{fig:qualitative}
\end{figure}

\subsection{Failure Modes}

\todo{Analyse and describe after evaluation. Expected failure cases based on
prior work and dataset properties:}

\begin{itemize}
    \item \textbf{Rare or abstract concepts} — icons for uncommon concepts
          (e.g., domain-specific technical symbols) that appear infrequently
          in training may produce generic or semantically misaligned outputs.

    \item \textbf{Complex multi-element compositions} — icons described as
          having more than 3--4 distinct visual elements tend to have lower
          rendering fidelity, consistent with the SVGenius finding that quality
          degrades with rising SVG complexity.

    \item \textbf{Invalid path data} — occasionally the model generates
          syntactically plausible but geometrically malformed \texttt{d}
          attributes (e.g., missing coordinate pairs after a command letter).
          This manifests as rendering errors. \todo{Quantify rate.}

    \item \textbf{ViewBox mismatch} — some generated icons place content
          outside the declared viewBox, resulting in clipped rendering.
\end{itemize}

\subsection{Caption Quality Ablation}

\todo{Compare: (1) VLM-generated captions vs. (2) raw icon names as prompts
for fine-tuning. This isolates the contribution of the captioning pipeline.}

\subsection{Comparison with General SVG Models}

\todo{Discuss qualitative differences between \method{} and StarVector/OmniSVG
on icon-specific prompts. Hypothesis: general models produce more complex,
less icon-like SVG structure; \method{} produces compact, \texttt{currentColor}
convention-following icons consistent with professional icon libraries.}
